\section*{Заключение}

В ходе выполнения работы были исследованы алгоритмы цифровой квадратурной демодуляции и многоскоростной фильтрации:
\begin{enumerate}
    \item Реализован перенос радиосигнала на нулевую промежуточную частоту с децимацией в $R=10$ раз с помощью каскада CIC-дециматора 8-го порядка и корректирующего КИХ-фильтра.
    \item Установлено, что динамический диапазон тракта (SFDR) при $nobg \ge 13$ бит ограничивается шумами квантования входного АЦП, поэтому избыточное наращивание разрядности гетеродина нецелесообразно.
    \item Показана критическая чувствительность тракта к квадратурному дисбалансу: амплитудная и фазовая асимметрия каналов приводят к резкому падению SFDR из-за роста неподавленной зеркальной составляющей.
    \item Эксперименты с ЛНА- и ЛЧМ-сигналами подтвердили линейность фазочастотных характеристик и сохранение формы широкополосных импульсов при каскадной обработке.
    \item Сравнение с однокаскадным КИХ-дециматором показало, что архитектура CIC + КИХ обеспечивает сопоставимую избирательность и подавление наложения спектров при существенно меньших затратах вычислительных ресурсов.
\end{enumerate}
